\chapter[Structural Risk and Graph Analysis]{Structural Risk and Graph Analysis}
\label{cp:network-sci}

% TODO: chapter not yet written. Decisions: docs/adr/0001; vocabulary: CONTEXT.md.
% Central claim: a component's position in the dependency graph carries risk information that
%   flat-inventory analysis cannot express; graph measures make it visible at three scales.
% Organise by scale, not by algorithm:
%   Component scale  -- articulation points and bridges, HITS hub/authority, spectral exposure score
%                       (eigenvector centrality; never "R0"), k-core number; one paragraph on the
%                       optimizer's edge rankings as an application of the spectral score.
%   Cluster scale    -- Louvain communities, bow-tie decomposition, rich-club.
%   Whole-graph scale-- spectral radius and epidemic threshold 1/lambda_1, degree entropy, cycles/SCCs,
%                       version conflicts (not "diamonds"), namespace concentration, power-law fit,
%                       small-world (C, L). Define briefly here; values are reported in Chapter 7 (RQ1).
% Out: blast radius and dominated sets (Chapter 5, refer back); unmaintained check (short note in
%   Chapter 5); WL fingerprinting / SBOM comparison (left out of the thesis); propagation simulation
%   (future work).

\section{Component Scale}
\label{sec:component-scale}

\subsection{HITS Hub and Authority Scores}
\label{sec:hits}

\section{Cluster Scale}
\label{sec:cluster-scale}

\section{Whole-Graph Scale}
\label{sec:structural-risk}
